\uuid{BdRn}
\titre{Raisonnements}
\auteur{bodin ; croizat ; sacre}
\organisation{exo7}
\licence{CC BY-NC-SA 4.0}
\datecreate{2022-03-08}
\niveau{L1}
\module{Algèbre}
\chapitre{Logique, ensemble, raisonnement}
\sousChapitre{Absurde et contraposée}
\difficulte{3}
\theme{Raisonnements}
\exo7qcmid{hilisit.logique.052}

\contenu{
\begin{question}
\qclassification{100.03}
Pour montrer une implication $\mathcal{P} \implies \mathcal{Q}$ par contraposition :
\begin{answers}
	\bad{Je suppose $\mathcal{P}$ et je montre $\mathcal{Q}$.}
	\bad{Je suppose $\mathcal{Q}$ et je montre $\mathcal{P}$.}
	\bad{Je suppose $\text{non}(\mathcal{P})$ et je montre $\text{non}(\mathcal{Q})$.}
	\good{Je suppose $\text{non}(\mathcal{Q})$ et je montre $\text{non}(\mathcal{P})$.}
\end{answers}
\begin{explanations}
Je suppose $\text{non}(\mathcal{Q})$ et je montre $\text{non}(\mathcal{P})$.
Par exemple pour montrer "$n^2$ pair $\implies$ $n$ pair", par contraposition je suppose $n$ non pair (c-à-d impair)
et je prouve alors que $n^2$ est aussi non pair.
\end{explanations}
\end{question}
}
